% ============================================================
%  Chapter 3: Biological Model
% ============================================================

\section{Биологическая модель}
\label{sec:biomodel}

\subsection{Перевод семи вопросов на язык биоинформатики}

\subsubsection{Вопрос 1: Определения ПВГ в биологии}

В биоинформатике существуют три канонических класса ПВГ.

\paragraph{Сети коэкспрессии генов (GCN).}
Пусть дана матрица экспрессии $X \in \mathbb{R}^{n \times m}$,
где $n$ — число генов, $m$ — число образцов.
ПВГ $G = (V, E, w)$ строится как:
\[
    V = \{\text{ген}_1, \ldots, \text{ген}_n\},
    \quad
    w(\{i, j\}) = \mathrm{cor}(X_i, X_j)
    = \frac{(X_i - \bar{X}_i)^\top (X_j - \bar{X}_j)}
           {\|X_i - \bar{X}_i\| \|X_j - \bar{X}_j\|},
\]
где $X_i \in \mathbb{R}^m$ — вектор экспрессии гена $i$.
Полученный граф полный ($K_n$), а веса принадлежат $[-1, 1]$~\cite{wgcna2008}.

\paragraph{Матрицы расстояний 3D-структур белков (Contact Maps).}
Пусть белок состоит из $n$ аминокислот с координатами $\mathbf{r}_1, \ldots, \mathbf{r}_n \in \mathbb{R}^3$.
Тогда:
\[
    w(\{i, j\}) = \|\mathbf{r}_i - \mathbf{r}_j\|_2, \quad w \in [0, +\infty).
\]
Сервер DALI~\cite{dali1993} сравнивает белки именно как ПВГ с метрическими весами.

\paragraph{Филогенетические сети.}
Узлы — геномы разных видов или штаммов.
Вес ребра $(i,j)$ — расстояние Хэмминга между геномными последовательностями:
\[
    w(\{i, j\}) = d_H(\mathbf{s}_i, \mathbf{s}_j)
    = \frac{1}{L}\sum_{k=1}^{L} \mathbf{1}[s_i^{(k)} \ne s_j^{(k)}],
\]
где $L$ — длина выровненных последовательностей.

\subsubsection{Вопрос 2: Практические задачи}

ПВГ-изоморфизм востребован в следующих задачах:
\begin{enumerate}
    \item \textbf{Поиск функции нового белка} (Protein Function Prediction):
          найти белок из базы данных PDB~\cite{alphafold2021},
          чья матрица расстояний изоморфна матрице нового белка.
    \item \textbf{Сравнительная транскриптомика}:
          установить, насколько GCN человека $\varepsilon$-изоморфна GCN мыши
          (обоснование переноса результатов с модельного организма).
    \item \textbf{Биомаркерный анализ} (Biomarker Discovery):
          найти под-граф GCN, где изоморфизм нарушен при онкогенезе.
\end{enumerate}

\subsubsection{Вопрос 3: Подходы к определению изоморфизма}

Точный изоморфизм в биологии практически недостижим из-за шума.
Применяются три класса релаксаций:

\paragraph{Глобальное выравнивание сетей (Global Network Alignment).}
Алгоритм IsoRank~\cite{isorank2008} решает задачу
$\max_{\varphi \in \mathfrak{S}_n} \sum_{i,j} W_1[i,j] \cdot W_2[\varphi(i),\varphi(j)]$,
что соответствует поиску биекции, максимизирующей совпадение весов рёбер.
Используется математика, аналогичная PageRank.

\paragraph{$\varepsilon$-изоморфизм (настоящая работа).}
Определение~\ref{def:eps_iso} и метрика $d_{\mathrm{iso}}$
обеспечивают строгую количественную меру несходства.

\paragraph{Оптимальный транспорт (Gromov–Wasserstein).}
Расстояние $d_{\mathrm{GW}}$ допускает нецелые «мягкие» соответствия,
устойчиво к шуму и вычисляется за полиномиальное время
(итеративная оптимизация)~\cite{gromov_wasserstein}.

\subsubsection{Вопрос 4: Резонанс с классами задач анализа данных}

Формализм ПВГ-изоморфизма резонирует со следующими задачами:
\begin{itemize}
    \item \textbf{Поиск аномалий} (Anomaly Detection):
          «сломанный» изоморфизм — это аномальный подграф.
    \item \textbf{Кластеризация} (Community Detection):
          функциональные модули генов суть клики в ПВГ с высокими внутренними весами.
    \item \textbf{Построение филогенетических деревьев}:
          минимальное остовное дерево ПВГ эволюционных расстояний (алгоритм Краскала).
    \item \textbf{Перенос обучения} (Domain Adaptation):
          две области данных изоморфны, если их граф-структуры метрически близки.
\end{itemize}

\subsection{Биологически-мотивированная версия теоремы консервативности}

Теорема~\ref{thm:conservation} формализует \emph{принцип Маркова–Хорвата}
в теории WGCNA~\cite{wgcna2008}:
\begin{quote}
\itshape
«Структурное сходство GCN двух организмов сильнее предсказывает функциональное
сходство их генов, чем сходство первичных последовательностей.»
\end{quote}

\subsection{Задача дифференциального анализа сетей}

\begin{definition}[Дифференциальный подграф]
Пусть $G_1 = \mathrm{GCN}(\text{норма})$ и $G_2 = \mathrm{GCN}(\text{рак})$.
\emph{Дифференциальным подграфом} на уровне $\theta$ называется граф
\[
    G_{\Delta} = \bigl(V,\; E_{\Delta},\; w_\Delta\bigr),
\]
где $E_{\Delta} = \{(i,j) : |W_2[i,j] - W_1[i,j]| > \theta\}$
и $w_\Delta(i,j) = |W_2[i,j] - W_1[i,j]|$.
\end{definition}

Вершины, инцидентные наибольшему числу рёбер в $G_\Delta$, являются кандидатами
в \emph{онкогенные маркеры}: гены, чья система взаимодействий радикально перестроилась.
Этот подход продемонстрировал высокую эффективность в практике~\cite{diff_network_cancer}.

\begin{example}[DALI и AlphaFold]
Система AlphaFold~\cite{alphafold2021} предсказала структуры $\sim 2 \times 10^8$ белков.
Для аннотации функции каждого нового белка $P_{\mathrm{new}}$ строится его матрица расстояний
$W_{\mathrm{new}}$, и ищется белок $P^*$ в базе данных PDB такой, что
\[
    P^* = \arg\min_{P \in \mathrm{PDB}} d_{\mathrm{GW}}(G(P_{\mathrm{new}}),\, G(P)).
\]
Сервис DALI реализует именно эту процедуру и является золотым стандартом
структурного выравнивания~\cite{dali1993}.
\end{example}
